\subsection{Implementation}
\label{subs:imp}

\myparagraph{\giptable}.
\giptable runs as a daemon process on each PU (Linux) to monitor local networking rules and coordinate with other PUs within the same computer.
It maintains ports and intercepts network rules within a network namespace.
\giptable leverages existing networking techniques, e.g., nsenter, NAT, policy routing, and iptables.
In \sys, we extend Kubernetes' CNI with \giptable.

\myparagraph{\gsocket.}
\gsocket is written in C/C++.
We leverage \codeword{LD\_PRELOAD}~\cite{ld-preload} to replace socket-related functions (implemented in libc) of an application at runtime.
Besides, we implement global shared memory (GShm) based on CXL (on CXL-DPU) and RDMA (on Bluefield-2).
The \gsocket implements a full set of Linux socket interfaces. 
We extend Linux kernel to support the speculative allocation and the interfaces of U-NAPI.

\myparagraph{Multi-ISA supports.}
We use the multi-platform container image feature~\cite{multi-platform} to support multi-ISA PUs in the same computer, which is a standard feature in Docker/Kubernetes.

\myparagraph{Security.}
\os is implemented as a per-process library now, which will maintain the global shared memory between two entites in userspace.
As each process can only access the shared memory for its own connections, \os will not introduce new security flaws to attack other processes.
A malicious process may break the connection, but can not cause other issues like overflow because of careful pointer operations.
Moreover, \os allows administrator to configure which connections are allowed to utilize the \gsocket by a configuration file, 
which is similiar like firmware to explictly specify the allowed list.



\myparagraph{Experimental setup.}
We use two settings for evaluation.
First (\textbf{Bluefield-DPU}), we use two servers, each with 2xE5-2650 v5 CPU (12cores, 2.2GHz each CPU), 96GB DDR4-1600 memory, and one Nvidia Bluefield-2 DPU devices.
Each Bluefield-2 DPU has 8 ARM Cortex-A72 cores (2.75GHz), 16GB DDR-1600 DRAM and ConnectX-6 (2x 100Gbps RDMA ports) NIC.
We also validate the results on Bluefield-3.
Second (\textbf{CXL-DPU}), 
we use a dual-socket server with Intel Xeon Platinum 8468V (2.9 GHz, 48 cores) and a 32GB 4800MT/s DDR5 RDIMM is connected through a CXL memory expansion card (CXL 2.x).
We implement CXL-based DPU with two VMs (representing CPU and DPU) and utilize the real CXL device (configured as \codeword{dax} dev) as the GShm.

